\documentclass{article}
\usepackage[utf8]{inputenc}
\usepackage[T1]{fontenc}
\usepackage{geometry}
\geometry{
    a4paper,
    left=2.5cm,
    right=2.5cm,
    top=2.5cm,
    bottom=2.5cm
}
\usepackage{graphicx}
\usepackage{float}
\usepackage{hyperref}
\usepackage[english]{babel}
\usepackage{amsmath}
\usepackage{amssymb}
\usepackage{array}
\usepackage{booktabs}
\usepackage{caption}
\usepackage{subcaption}
\usepackage{enumitem}
\usepackage{fancyhdr}
\usepackage{titlesec}
\usepackage{url}
\usepackage{textcomp}
\usepackage{xcolor}
\usepackage{listings}
\usepackage{multirow}
\usepackage{longtable}
\usepackage{tabularx}
\usepackage{times}

% Set up page style
\pagestyle{fancy}
\fancyhf{}
\rfoot{\thepage}
\renewcommand{\headrulewidth}{0pt}

% Title formatting
\titleformat{\chapter}[display]{\normalfont\bfseries}{\chaptertitle\thechapter}{0pt}{\Huge}
\titlespacing*{\chapter}{0pt}{50pt}{40pt}

% Document info
\title{Template Document}
\author{Academic Support Agent}
\date{\today}

\begin{document}
\maketitle

\doublespacing
\begin{titlepage}


\setlength{\droptitle}{-4em} 
\center
{\huge \bfseries PROJECT REPORT } % Title of the project
\\[1.5cm]
{\huge \bfseries STUDENT ATTENDANCE MANAGER APP } % Title of the project
\\[1cm] \begin{minipage}{0.7\textwidth}
\begin{center} \large
\textbf {Course Code: CSIG 364}\\[0.5cm]
\textbf {Course Name: Mini Project Using Java}\\[0.5cm]

\end{center}
\end{minipage} 

%\addtolength{\droptitle}{-4pt}
\vspace{10em}

\includegraphics[width=4cm]{UNILOGO.jpg}

\vspace{10em}
\begin{minipage}{0.5\textwidth}
\begin{center} \large
\textbf {Submitted By: MITRA PRANOY PANDI} % Name of the student
\\[0.5cm] 
\textbf {Registration no: 20384118} \\[0.5cm] % Registration number of the student
\textbf{Department of Computer Science} % Name of the department
\\[0.5cm]
\textbf{Mini Project} % Type of meeting
\end{center}
\end{minipage} 

\vspace{3em}
{\large \today}
\end{titlepage}
%contents


%certificate
\newcommand{\RomanNumeralCaps}[1]
    {\MakeUppercase{\romannumeral #1}}
\begin{center}
\thispagestyle{empty}

\LARGE{\textbf{Pondicherry University}} \\ 
\large{\textbf{Department of Computer Science}}\\
\large{\textbf{Kalapet,Pondicherry-605014}}\\[0.5cm]

\includegraphics[scale=0.5]{UNILOGO.jpg}\\[0.5cm]

{\Huge \textbf{BONAFIDE CERTIFICATE}}\\[0.5cm]
\end{center}
\linespread{1.13}
\large{This is certify that the Mini Project entitled 
\textbf{`` Student Attendance Manager App'',} 
submitted by 
\textbf{Mitra Pranoy Pandi [Regno: 20384118]}
 is a record of bonafide work carried out by him, in the partial
 fulfilment of the requirement for the award of Integrated Msc.Computer Science at Pondicherry University This work is done in
  \RomanNumeralCaps{6} sem during  year 2022-2023, under  guidance.}\\[0.5cm]
\begin{flushright}
    \large{\underline{\textbf{Project Guide}}}\\[0.5cm]
    \textbf{{ Dr. Sukhvinder Singh}}\\[0.1cm]
    \textbf{{ Asst. Professor }}\\[0.1cm]
    \textbf{{Dept.Computer Science \& Engineering}}\\
\end{flushright}

\begin{flushleft}

    \Large{\underline{\textbf{External}}}\\[0.8cm]
\large{\textbf{Examiner Name:}}\\[0.8cm]
\large{\textbf{Signature}}\\[0.8cm]
\large{\textbf{Date:}}
\end{flushleft}
\begin{flushright}
    \large{\textbf{Guide Signature}}\\[0.5cm]
\end{flushright}

%endofcertificate
\newpage
\doublespacing
\section*{\centering{Acknowledgement}}
It gives me a great pleasure in presenting the mini project report on 'ATTENDANCE MANAGEMENT APP'. I would like to express my deepest gratitude to my faculty,  Dr. Sukhvinder Singh his valuable guidance, consistent encouragement for doing this project.  I am really grateful to them for their support and giving me all the help and guidance I needed. Their valuable suggestions were very helpful.\\
\begin{flushright}
{\bf Mitra Pranoy Pandi}\\[0.2cm]
{\bf 20384118}
\end{flushright}
\newpage

\newpage
\doublespacing
\section {Objective Of Course}
{\fontsize{14}{14}\selectfont The objective of a course on "mini project in Java" would typically be to provide students with an opportunity to apply the concepts and skills they have learned in Java programming to develop a small-scale software project.\\
Through this course, students would learn about project management, software development life cycle, software design patterns, database design, and user interface design and also gain experience in working with Java development tools and technologies, such as \textbf{Intellij , Android Studio , Java FX.}
\\
The course would typically involve the following activities:
\begin{enumerate}
   \item Defining the project scope and requirements
  \item Planning and scheduling the project
  \item Designing the software architecture and user interface
  \item Implementing the software using Java programming language
  \item Testing and debugging the software
  \item Documenting the project
\end{enumerate}
}
\newpage
\doublespacing
\section{Abstract}
%begining of abstract
{\fontsize{14}{14}\selectfont The Attendance Manager Android App is a mobile application designed to help teachers and professors manage attendance records of their students. \\
This app is user-friendly, efficient, and can be accessed from anywhere, anytime. The app enables teachers to add, edit, and delete courses and students, record attendance, and generate reports. \\
The app has been developed using Java programming language, Android Studio IDE, and Realm database. The app meets all the functional and non-functional requirements, and has been tested for compatibility with Android OS version 7.1 and above. \\
The app is an essential tool for teachers and professors to manage their attendance records, reduce manual effort, and improve the efficiency of their work.}

\newpage
\section{Problem Objective}
\doublespacing
%problem objective
{\fontsize{14}{14}\selectfont The Problem objective of an attendance manager is to efficiently and accurately track the attendance of students. This involves creating a system for recording attendance, maintaining attendance records, generating reports, and enforcing attendance policies.\\

The primary goal of an attendance manager is to ensure that attendance policies are being followed and that students are attending classes as required. This helps organizations to identify any attendance-related issues, such as excessive absences or tardiness, and take appropriate action to address them.\\

In addition, an attendance manager may use attendance data to identify patterns and trends, such as peak absenteeism times or departments with higher rates of absenteeism, which can help inform organizational policies and practices.\\

Overall, the objective of an attendance manager is to promote a culture of attendance and accountability within the organization or educational institution.}

\newpage
\section{Problem Statement}
\doublespacing
%problem statement
{\fontsize{14}{14}\selectfont The problem statement for an attendance manager app can be defined as the need for an efficient and user-friendly tool to manage attendance records for organizations or educational institutions.\\

Manual attendance tracking can be time-consuming and prone to errors, leading to inaccurate records and potential compliance issues. Additionally, with the rise of remote work and online learning, traditional methods of attendance tracking may not be feasible or effective.\\

Therefore, there is a need for a digital solution that allows for easy and accurate attendance tracking, as well as providing real-time visibility into attendance data. The attendance manager app should be designed to simplify the process of recording attendance, generating reports, and enforcing attendance policies.\\

Overall, the problem statement for an attendance manager app is the need for a reliable, efficient, and user-friendly tool to manage attendance records and ensure compliance in a modern work or educational environment}

\newpage
\section{Function}
\doublespacing
%function
{\fontsize{14}{14}\selectfont
The following are the functional requirements of the Attendance Manager Android App:
\begin{enumerate}
\item Add, edit courses
\item Add, edit,students in each course
\item Record attendance for each student in each course
\item Generate attendance reports for each course and student
\item View attendance records of all students in all courses
\end{enumerate}}

\newpage
\section{Design}
\subsection{UML Diagrams}
\doublespacing
{\fontsize{14}{14}\selectfont
UML diagrams are graphical representations used in software engineering to model, visualize, and document software systems. They help to communicate the design of a system and provide a way to depict the system's architecture, behavior, and interactions between its components.\\

UML diagrams can be classified into three categories: structural diagrams, behavioral diagrams, and interaction diagrams. Examples of UML diagrams include use case diagrams, class diagrams, activity diagrams, sequence diagrams, and more. By using UML diagrams, software developers can better understand and communicate the design of their software systems, making it easier to collaborate with stakeholders and ensure that the system meets their requirements.}
\subsubsection{Use Cases Diagram}
\includegraphics[width=18cm,height=20cm]{UML diagrams/UseCaseDiagram.jpg}
\subsubsection{Activity Diagram}
\includegraphics[width=18cm,height=20cm]{UML diagrams/Attendance Manager ActivityDiagram.jpg}
\subsubsection{Class Diagram}
\includegraphics[width=18cm,height=22cm]{UML diagrams/Attendace manager ClassDiagram.jpg}
\subsubsection{ER-Diagram}
\includegraphics[width=18cm,height=15cm]{UML diagrams/ERDDiagram.jpg}

\newpage
\subsection{Project Structure}
\begin{figure}[htbp]
 \includegraphics[width=8cm]{UML diagrams/Project Struct.png} \hfill
 \includegraphics[width=7cm]{UML diagrams/Projectsruct2.png}
 \caption{Project Structure In Android Studio}
\end{figure}
\section{Sample Screens}
\includegraphics[width=15cm,page=1]{UML diagrams/Sample Screns.pdf}
\newpage
\includegraphics[width=15cm,page=2]{UML diagrams/Sample Screns.pdf}
\newpage
\includegraphics[width=15cm,page=3]{UML diagrams/Sample Screns.pdf}
\doublespacing
\section{Methodology}
{\fontsize{14}{14}\selectfont
Developing a methodology for  mini project attendance manager app  involve the following steps:
\begin{enumerate}

   \item \textbf{Requirements:} A Laptop with minimum configuration of 8GB Ram ,256GB Storage and i5 or ryzen 3500u or higher,Android Studio, Java Programming Knowledge.

\item \textbf{Define project scope and objectives:} Begin by identifying the key objectives of the attendance manager app. Determine the specific requirements that need to be fulfilled, such as the ability to capture attendance data, generate reports. Determine the scope of project by identifying the functionalities that plan to include in  app.

 \item \textbf{Select a development methodology:}  select a development methodology to guide the development process. A popular methodology for small projects is Agile. The Agile methodology emphasizes flexibility and iterative development. I used Agile methodology in this project

 \item \textbf{Plan and design the app:} Plan the app's features, user interface, and database schema. Sketch out the app's UML and create a visual design that is simple and user-friendly.

 \item \textbf{Develop the app:} Begin the actual development process by writing code, integrating the database, and adding the features. Ensure that the app is user-friendly, secure, and scalable. 
 
 
\end{enumerate}}
\newpage
\section{Applications Of Project}
%applications
{\fontsize{14}{14}\selectfont
An attendance manager app can be useful in a variety of settings, including:
\begin{enumerate}
   \item Schools and universities: Attendance management is an important part of tracking student progress and ensuring that they are meeting academic requirements. An attendance manager app can be used by teachers to record attendance, track absences, and generate reports.

  \item Businesses and organizations: Attendance management is essential in workplaces to ensure that employees are present and productive. An attendance manager app can help managers track employee attendance, manage schedules, and generate payroll reports.

  \item Events and conferences: Attendance management is crucial for organizing events and conferences. An attendance manager app can help organizers keep track of attendees, manage registration and check-in processes, and generate reports.

  \item Fitness centers and gyms: Attendance management is important in fitness centers and gyms to ensure that members are attending regularly and meeting their fitness goals. An attendance manager app can help gym managers track attendance, manage schedules, and monitor progress.

  \item Religious organizations: Attendance management is important for religious organizations to track attendance at services and events. An attendance manager app can help religious leaders keep track of attendance, manage schedules, and generate reports.

Overall, an attendance manager app can be useful in any setting where attendance tracking is required. It can save time, reduce errors, and improve efficiency in managing attendance.
\end{enumerate}}
\newpage
\section{References}
{\fontsize{14}{14}\selectfont
Here are the references I used in creating my project and not everything is directly used in the project changes are made according to my comfort.
\begin{enumerate}
   

\item  Android Studio Developer Docs For developing apps
\url{https://developer.android.com/guide}
\item Realm Database which is developed by MongoDb Docs
\url{https://www.mongodb.com/docs/realm/}

\item For Debugging and Writing the code I refer Stack Overflow ,Github and Stack Exchange.
\url{https://stackoverflow.com/},\\\url{https://github.com/},\\\url{https://android.stackexchange.com/}
\end{enumerate}
\end{document}